\begin{resumo}[Abstract]
\begin{otherlanguage*}{english}
Interrupt-driven embedded software suffers from concurrency defects where an interrupt service routine may fire between two
accesses to a shared variable and corrupt its value. Existing tools favour
precision, measure recall poorly, and only the oldest one proposes fixes.
Meanwhile, recent work uses large language models to triage and repair
static-analyzer alerts, but all of it addresses sequential defects. This work
sits at the empty intersection between the two literatures. A static analysis
over LLVM's intermediate representation finds candidates with no false
negatives relative to an explicit list of assumptions, an SMT solver discards
only what it proves impossible, and a language model ranks, explains and
proposes fixes for the remaining candidates without ever discarding them. The
evaluation uses the Racebench benchmark with a recounted and corrected ground
truth, metrics that cannot be improved by discarding candidates, namely the
recall gate, trap rejection and the Inspection Ratio, and an ablation of prompt
techniques. This partial report presents the conceptual background, the work
method, the requirements, the tool design and the evaluation procedure.

\vspace{\onelineskip}

\noindent
\textbf{Keywords}: concurrency. interrupts. embedded software. static analysis.
large language models.
\end{otherlanguage*}
\end{resumo}
